\scanpage{179}
% archive sheet number: 178
tels que \deleted{\unsure{le}} $f_i$ soit un coéfficient pour $V_i$, alors
il suffit de prendre $V=\sum V_i$. \deleted{\unsure{Nécessaire, car}} \inserted{Suffisant,}
car considérant \deleted{tous les coéfficients de la} \inserted{une base $(e_i)$} de $V$,
et les coéfficients $f_{e_i,e'_j}$, on voit que
tout coéfficient relatif à $\overset{\unsure{n}}{\otimes}V$ (donc à
un sous-truc ou un truc quotient) est
contenu dans l'algèbre engendrée par les
$f_{e_i,e'_i}$.
% slanted
\margin{par définition ou \unsure{paraphrase} de la multiplication dans la coalgèbre $A_0$, on a, pour \unsure{des} $A$-comodules $V,W$ et $x\,x'\,y\,y'$ dans $V\,V'\,W\,W'$,\[\begin{aligned}f_{x,x'}f_{y,y'}&=f_{x\otimes y,x'\otimes y'}.\ \text{\unsure{V.}}\\ f_{x,x'}\unsure{*}f_{y,y'}&=f_{x\otimes x',y\otimes y'}\ \text{\unsure{''}}\end{aligned}\]}
% slanted
\margin{\deleted{\illegible{1}} \uline{Lemme} 1. L'algèbre engendrée par les coefficients de la repr. dans $V$ est l'ensemble des coefficients dans \deleted{\unsure{tous}} les représentations $\overset{n}{\otimes} V$. De plus, les représentations \unsure{contenues} dans un \unsure{quotient} d'un $\overset n\otimes V$ forment une \uline{classe} de représentations.}

\uline{N.B.} \deleted{Soit} Si $G$ est algébrique linéaire
il a une représentation linéaire fidèle \deleted{\unsure{pour}}
\deleted{\inserted{\unsure{d'image fermée}}}, $\searrow\ \oplus$ \deleted{\unsure{induisant la structure ci-dessus, la représentation dans $V$.}}
\inserted{qui est un \uline{isomorphisme} $\updownarrow$ \unsure{onto.} au sens de l'énoncé des conditions équivalentes \unsure{barrées}.}
i.e. $G$ est
un sous-groupe ``fermé'' de $\mathrm{Gl}(V)$, avec
la structure algébrique induite (i.e. $A_0$ est
l'algèbre des polynômes en les fonctions
``coordonnées'' $g_{ij}$ et $\det(g_{ij})^{-1}$) : il
suffit de prendre \deleted{un \unsure{système (f) fini}}
\deleted{de générateurs de $A_0$, et une} $V$
comme ci-dessus, i.e. telle que $A_0$ soit
l'alg. engendrée par les \deleted{\unsure{fonctions}} $g_{ij}\in A_0$.
% slanted
\margin{en fait, il faut dire que \deleted{\unsure{toutes}} représentations de $R$ sont \uline{\unsure{exactt|exactement}} les représentations ``régulières'' du groupe algébrique $G$ défini de la façon voulue}
Il faut seulement montrer a) que $\det(g_{ij})^{-1}\in A_{\unsure{0}}$
ce qui résulte du fait que $\det(g_{ij})$ ne
s'annule pas sur $G$ donc \deleted{\unsure{est un inv.}} \inserted{son inverse}
est \deleted{\unsure{rationnelle}} régulière sur $G$ ; b) que
$\overline G\subset\mathrm{Gl}(W)$ est fermé pour la top.
de Zariski de $\mathrm{Gl}(W)$ (et $=$ dans $L(W)$,
pour la top de Zariski de $L(\unsure{V|W})$). Mais
